\documentclass[letterpaper,11pt]{report}
\usepackage{color}
\usepackage{fullpage}
\usepackage{listings}
\usepackage{caption}
\usepackage[bookmarks=true]{hyperref}
\usepackage{graphicx}

\usepackage{syntax}
\setlength{\grammarparsep}{20pt plus 1pt minus 1pt} % increase separation between rules
\setlength{\grammarindent}{5em} % increase separation between LHS/RHS
\AtBeginDocument{\catcode `\_=8}

\usepackage{fouriernc}
\usepackage[T1]{fontenc}

%\usepackage[scaled]{beramono}
%\usepackage[T1]{fontenc}

%\usepackage{inconsolata}
%\usepackage[T1]{fontenc}

\title{bawk, a binary awk}
\author{
	Kevin M.\ Graney\\
	\texttt{kmg2165@columbia.edu}
}

\lstset{ %
linewidth=\textwidth,
basicstyle=\footnotesize\renewcommand{\ttdefault}{pcr}\ttfamily,       % the size of the fonts that are used for the code
numbers=left,                   % where to put the line-numbers
numberstyle=\tiny\sf,      % the size of the fonts that are used for the line-numbers
stepnumber=1,                   % the step between two line-numbers. If it's 1 each line 
                                % will be numbered
numbersep=5pt,                  % how far the line-numbers are from the code
backgroundcolor=\color{white},  % choose the background color. You must add \RequirePackage{color}
fillcolor=\color{white},
showspaces=false,               % show spaces adding particular underscores
showstringspaces=false,         % underline spaces within strings
showtabs=false,                 % show tabs within strings adding particular underscores
frame=trlb,
tabsize=4,	                  % sets default tabsize to 2 spaces
captionpos=t,                   % sets the caption-position to top
belowcaptionskip=2pt,
breaklines=true,                % sets automatic line breaking
breakatwhitespace=false,        % sets if automatic breaks should only happen at whitespace
title={\texttt{\lstname}},                 % show the filename of files included with \lstinputlisting;
upquote=true,
numberbychapter=true,
escapeinside={\%*}{*)},         % if you want to add a comment within your code
morekeywords={*,...},           % if you want to add more keywords to the set
keywordstyle=\color[rgb]{0,0,1},
commentstyle=\color[rgb]{0.133,0.545,0.133},
stringstyle=\color[rgb]{0.627,0.126,0.941},
%keywordstyle=\bfseries,
%commentstyle=\itshape,
%stringstyle=\slshape,
frameround=fttf,
framexleftmargin=14pt,
xleftmargin=14pt,
xrightmargin=4pt,
language=awk
}

\DeclareCaptionFont{white}{\color{white}}
\DeclareCaptionFormat{listing}{\setlength{\fboxsep}{2pt}\hskip-4pt\colorbox{black}{\parbox{470pt}{#1#2#3}}}
\captionsetup[lstlisting]{format=listing,labelfont=white,textfont=white, singlelinecheck=false, margin=0pt, font={bf,footnotesize}}

\newcommand{\testfile}[1]{%
	\subsection{#1}
	\lstinputlisting[caption=Test program (#1),language=awk]{../tests/#1}
	\lstinputlisting[caption=Expected output (#1.out),language=awk]{../tests/#1.out}}

\begin{document}
%\maketitle
\begin{titlepage}
\begin{center}
~\\\vspace{1.75in}
\Huge
\textsc{bawk, a binary awk}

\vspace{1in}
\LARGE
Kevin M.\ Graney\\
\texttt{kmg2165@columbia.edu}

\Large
\vspace{1.75in}
Programming Languages and Translators\\
COMS W4115, Spring 2013, Professor Stephen Edwards

\end{center}
\vfill
\includegraphics[width=4in]{columbia_university.eps}
\end{titlepage}
\tableofcontents
\lstlistoflistings

\chapter{Introduction}
\section{Overview}
The bawk language, who's name is derived from `binary awk', is intended to be a small, special-purpose language for the parsing of formatted binary files.  In the spirit of awk, bawk will match rules in binary files and extract data based on these rules.  Bawk will solve the problem of quickly decoding a binary file of a known format to extract information much the way awk solves this same problem for formatted text files.

The bawk interpreter will take two inputs for execution: (1) text in the bawk langauge, and (2) binary data to run the bawk program on.  This workflow will be similar to that of a basic awk workflow, and the language itself is inspired by awk, with the same basic program construct.  That is, a program will be a series of patterns and statements, with the statements executing when their respective pattern matches.

\section{Sample program}
A representative bawk program is shown below.  This program operates on PNG format image files and will print two numbers, newline separated.  The first number is the width of the image and the second number is the height. 
\lstinputlisting[caption=A simple bawk program]{../tests/test-useful1.bawk}


\chapter{Language tutorial}
\chapter{Language reference manual}
\section{Lexical Conventions}
A program consists of a single bawk language character file.  There are six classes of tokens: identifiers, keywords, constants, string literals, operators, and separators.  White space characters are ignored except as they separate tokens or when they appear in string literals.

\subsection{Comments}
Comments follow the ANSI C style, beginning with \texttt{/*} and ending with \texttt{*/}.  They do not nest, and they cannot be present in quoted strings.

\subsection{Identifiers}
\label{sec:identifiers}
An identifier is a sequence of uppercase and lowercase letters A--Z, the numerals 0--9, and the underscore character.  The first character of an identifier can not be a numeral.

\begin{grammar}
<identifier> ::= <letter>
\alt <identifier> (<letter> | <digit>)
\end{grammar}

\subsection{Keywords}
The following identifiers are reserved for use as keywords, and may not be used otherwise: \texttt{if}, \texttt{else}, \texttt{return}, \texttt{def}, \texttt{int1}, \texttt{int2}, \texttt{int4}, \texttt{uint1}, \texttt{uint2}, \texttt{uint4}.

\subsection{Constants}
\label{sec:constants}
Bawk only supports a single kind of constant: the integer constant.  The semantic meaning of integer constants depends on the context in which they occur.  Except for inside of pattern expressions (see \S\ref{sec:pattern-constants}), the following is true.  An integer constant consisting of a sequence of digits is taken to be in decimal.  A sequence of hexadecimal digits prefaced by \texttt{0x} is taken to be in hexadecimal (base 16).

\subsection{String Literals}
\label{sec:strings}
A string is a sequence of characters enclosed in double quotes.  The only valid place for strings is inside of pattern expressions (see \S\ref{sec:pattern-strings}).

\section{Meaning of Identifiers}
Identifiers can have several different meanings based on how they are first used.  Once an identifier is associated with a particular meaning it cannot be disassociated with that meaning.

\subsection{Function Name}
Previously unused identifiers become function names once a function is declared with a given name.

\subsection{Variable Name}
\label{sec:variables}
Previously unused identifiers become variable names once an assignment expression (see \S\ref{sec:assignment}) executes with the identifier on the left hand side.

\subsection{Pattern Binding Variable Name}
\label{sec:bind-variables}
Previously unused identifiers become binded variable names once a pattern expression with a binding pattern token (see \S\ref{sec:pattern-bindings}) matches.

\subsection{Special Identifiers}
\label{sec:special-identifiers}
There are certain special identifiers that behave like variables, but whose meaning affects the program execution.
\begin{itemize}
\item \texttt{RP} -- The \textsl{file pointer} variable  indicates where in the file pattern statements begin matching (see \S\ref{sec:pattern-statement}).  The value contains the number of bytes from the beginning of the file.  At the start of the bawk program this value defaults to $0$.
\item \texttt{LE} -- Force \textsl{little-endian} interpretation of integers and strings while pattern matching in the file.  A value of $0$ indicates big-endian interpretation and a value of $1$ indicates little-endian interpretation.  The default value is $0$.
\item \texttt{\_} -- The \texttt{trash} variable is a write-only identifier that cannot be read.  It is useful for binding don't-care values in pattern matching expressions.
\end{itemize}

\section{Statements}
\label{sec:statements}
The expressions described in \S\ref{sec:expressions} are a specific form of statement in the bawk language.  There are three different types of statements, and a list of statements forms a bawk program.  The three different types of statements are: expressions (described in \S\ref{sec:expressions}), block statements, and pattern statements.
\begin{grammar}
<program> ::= <statement-list>

<statement-list> ::= $\epsilon$
\alt <statement>
\alt <statement-list> <statement>

<statement> ::= <expression-statement>
\alt <block-statement>
\alt <pattern-statement>
\alt <function-declaration>
\end{grammar}

\subsection{Block Statements}
Block statements are used to combine a sequence of statements into a single statement.  The sequence of statements inside the block is executed in order when the block itself is executed.  Variables first used inside the block are locally scoped, but names from outside the block are also available.
\begin{grammar}
<block-statement> ::= `{' <statement-list> `}'
\end{grammar}

\subsection{Pattern Statements}
\label{sec:pattern-statement}
Pattern statements allow a type of pattern matching to be performed on the binary file.  Pattern statements consist of a pattern expression (see \S\ref{sec:pattern-expressions}) and a statement.  If the pattern expression matches the data file at the location of \texttt{RP} at the time of evaluation of the pattern statement, then the corresponding statement is executed.  If the pattern does not match, then the statement is not executed.
\begin{grammar}
<pattern-statement> ::= `/' <pattern-expression> `/' <statement>
\end{grammar}

\subsection{Expression Statements}
An expression statement is an expression followed by a semicolon.  Expressions are described in \S\ref{sec:expressions}.

\begin{grammar}
<expression-statement> ::= <expression> `;'
\end{grammar}

\subsection{Function Declaration}
A function declaration is a statement declaring a new function.  Function declarations consist of an identifer to serve as the function name, followed by an open parenthesis, an optional comma separated list of identifiers to serve as function parameter names, a closing parenthesis, and a statement to serve as the function body.

\begin{grammar}
<function-declaration> ::= <identifier> `(' <function-decl-params> `)' <statement> 

<function-decl-params> ::= $\epsilon$
\alt <identifier>
\alt <function-decl-params> `,' <identifier>
\end{grammar}

\section{Expressions}
\label{sec:expressions}

\subsection{Primary Expression}

\begin{grammar}
<primary-expression> ::= <constant>
\alt <identifier>
\alt <string-literal>
\alt <function-call>
\alt `(' <expression> `)'

<expression> ::= <assignment-expression>
\alt <additive-expression>
\end{grammar}

\subsection{Function Calls}
A function call is an identifier, known as the function name, followed by a pair of parenthesis containing a, possibly empty, comma-separated list of expressions.  These expressions constitute arguments to the function.  When functions are called a copy of each argument is made and used within the function, that is, all function calls pass arguments by value.
\begin{grammar}
<function-call> ::= <identifier> `(' <function-call-params> `)'

<function-call-params> ::= $\epsilon$
\alt <expression>
\alt <function-call-params> `,' <expression>
\end{grammar}

\subsection{Multiplicative Operators}
The multiplicative operators, \texttt{*} and \texttt{/}, are left associative.  The result of \texttt{*} is multiplication of the two operands, and the result of \texttt{/} is the quotient of the two operands.  The quotient with a second operator of zero yields an undefined result.
\begin{grammar}
<multiplicative-expression> ::= <primary-expression>
\alt <multiplicative-expression> `*' <primary-expression>
\alt <multiplicative-expression> `/' <primary-expression>
\end{grammar}

\subsection{Additive Operators}
The two additive operators, \texttt{+} and \texttt{-}, are left associative.  The expected arithmetic operation is performed on integers.  The result of the \texttt{+} operator is the sum of the operands, and the result of the \texttt{-} operator is the difference of the operands.
\begin{grammar}
<additive-expression> ::= <multiplicative-expression>
\alt <additive-expression> `+' <multiplicative-expression>
\alt <additive-expression> `-' <multiplicative-expresssion>
\end{grammar}

% Shift Operators?

\subsection{Relational Operators}
The relational operators are all left associative.  The operators all produce $0$ if the specified relation is false and $1$ if the specified relation is true.  The operators are \texttt{<} (less than), \texttt{>} (greater than), \texttt{<=} (less than or equal to), \texttt{>=} (greater than or equal to).
\begin{grammar}
<relational-expression> ::= <additive-expression>
\alt <relational-expression> `<' <additive-expression>
\alt <relational-expression> `>' <additive-expression>
\alt <relational-expression> `<=' <additive-expression>
\alt <relational-expression> `>=' <additive-expression>
\end{grammar}

\subsection{Equality Operators}
The \texttt{==} (equal to) and \texttt{!=} (not equal to) operators compare operands for equality.  Like the relational operators, the equality operators produce $0$ if the relationship is false and $1$ if the relationship is true.  The equality operators have a lower precedence than the relational operators, and are also left associative.
\begin{grammar}
<equality-expression> ::= <relational-expression>
\alt <equality-expression> `==' <relational-expression>
\alt <equality-expression> `!=' <relational-expression>
\end{grammar}

% Logical AND and Logical OR?

\subsection{Assignment Expressions}
\label{sec:assignment}
All assignment expressions require a modifiable identifier on the left hand side.  The only type of modifiable identifier is a variable (see \S\ref{sec:variables}).  The assignment operator, \texttt{=}, is right-associative and the return value of assignment is the value being assigned.
\begin{grammar}
<assignment-expression> ::= <equality-expression>
\alt <identifier> `=' <assignment-expression>
\end{grammar}


\section{Pattern Expressions}
\label{sec:pattern-expressions}
A pattern expression matches content in the binary data file.  Pattern expressions are distinct in syntax from the expressions described in \S\ref{sec:expressions}.

\begin{grammar}
<pattern-expression> ::= $\epsilon$
\alt <pattern-token>
\alt <pattern-expression> <pattern-token>

<pattern-token> ::= <pattern-constant>
\alt <pattern-binding>
\alt <pattern-variable>
\alt <string-literal>
\end{grammar}

Pattern expressions can contain constants or typed variable bindings.  Each constant or binding is known as a pattern term, and pattern terms are whitespace delimited.  Already bound pattern variables can also be used in pattern statements.

\subsection{Constants}
\label{sec:pattern-constants}
Inside pattern expressions, constants work differently than those described in \S\ref{sec:constants}.  All constants in pattern terms are implicitly expressed in hexadecimal form, and the \texttt{0x} prefix required to express hexadecimal constants in \S\ref{sec:constants} must be ommitted.

Constants are read in hexadecimal as bytes.  For this reason leading zeros on a value have semantic meaning.  For example \texttt{/0000abcd/} matches the four byte pattern \texttt{00 00 ab cd} in the binary file, while \texttt{/abcd/} matches the two byte pattern \texttt{ab cd}.

The byte order expected in the file is dependent on the run-time value of the \texttt{LE} special variable (see \S\ref{sec:special-identifiers}).  If this variable is non-zero, then constants are interpreted in \textsl{little endian} order.  If the variable is zero, then the interpretation is \textsl{big endian}.  The default value is zero.  As an example, if \texttt{LE != 0} the pattern \texttt{/aabb/} matches the file with two bytes in the order \texttt{bb aa}, and if \texttt{LE == 0} then the pattern matches bytes \texttt{aa bb} in that order.

Whitespace does make a difference in the interpretation of pattern constants.  When \texttt{LE != 0}, the patterns \texttt{/aabb/} and \texttt{/aa bb/} are distinctly different.  If \texttt{LE == 0}, these two patterns are the same.  Whitespace added between byte boundaries also changes the semantic meaning; \texttt{/01ab/} is \emph{not} identical to \texttt{/0 1 a b/}.  If a constant has an odd number of hexadecimal digits, an implied \texttt{0} is added to the left-most side of the value.  That is \texttt{/123/} is the same as \texttt{/0123/} (is the same as \texttt{/01 23/} if \texttt{LE == 0}).  For this reason the example of \texttt{/0 1 a b/} is identical to \texttt{/00 01 0a 0b/} (and \texttt{/00010a0b/} if \texttt{LE == 0}).
 
\subsection{Bindings}
\label{sec:pattern-bindings}
A bind pattern term consists of an identifier, a colon (\texttt{:}), and a bind type.  Valid bind types are: \texttt{int1}, \texttt{int2}, \texttt{int4}, \texttt{uint2}, \texttt{uint4}.  The \texttt{int1} type matches a one-byte integer, the \texttt{uint4} type matches an unsigned 4-byte integer.  Identifiers have the same naming rules as \S\ref{sec:identifiers}.
\begin{grammar}
<pattern-binding> ::= <identifier> `:' <bind-type>

<bind-type> ::= `int1'
\alt `int2'
\alt `int4'
\alt `uint1'
\alt `uint2'
\alt `uint4'
\end{grammar}

\subsection{Pattern Variables}
Previously bound pattern variables can be used in a pattern expression to match the previously bound value.
\begin{grammar}
<pattern-variable> ::= <identifier>
\end{grammar}

\subsection{String Literals}
\label{sec:pattern-strings}
String literals are also supported in pattern expressions.  The syntax is exactly the same as strings in \S\ref{sec:strings}.  The string literal is converted to and interpreted as a sequence of ASCII-encoded bytes, exactly how constants are (see \S\ref{sec:pattern-constants}).  





\chapter{Project plan}
\chapter{Architectual design}
\section{Design decisions}
The design decisions for this project were not made with considerations towards performance or scalability.  Design decisions were based around ease of implementation and demonstration of compiler concepts.

\section{Execution phases}
The bawk interpreter consists of a bytecode compiler and a bytecode runtime environment.  The two are seamlessly integrated into a single executable program, \texttt{bawk}.  By calling this program with a \texttt{-c} option the results of the bytecode compiler are output, and by calling it with no option (or an explicit \texttt{-e}) the the given program is compiled internally to bytecode and then immediately executed in the bytecode interpreter.

\subsection{Bytecode compiler}
\label{sec:bytecode-compiler}
The bytecode compiler consists of a scanner (generated by OCamllex), a parser (generated by OCamlyacc), and the compiler itself.  The scanner parses a sequence of input characters that form a given bawk language program and returns a list of tokens.  The token types include punctuation that has meaning in the language, the C-style comments (\texttt{/*...*/}), identifiers, keywords, constants, and string literals.  Whitespace is consumed by the scanner, but serves no purpose other than to separate tokens.

There is at least one instance of ambiguity in the scanner.  Identifier names, as defined in \S\ref{sec:identifiers}, consist of a letter or underscore followed by letters, underscores, and/or numerals.  However, constants in pattern expressions, as defined in \S\ref{sec:pattern-constants}, consist of a sequence of hexadecimal digits, possibly including the letters 'a'-'f' in uppercase or lowercase.  This leads to possible confusion in the code chunk listed below.  On line 2 it is not known if \texttt{cafe123} is a variable name or a hexadecimal number.  This ambiguity is actually resolved by the parser, which treats ambiguous constants as constants if they are valid hexadecimal.

\begin{lstlisting}[caption=Possibly ambiguous code]
/cafe123:int1/ {
	/cafe123/ {
		/* do_something */
	}
}
\end{lstlisting}

The parser converts the sequence of tokens into an abstract syntax tree (AST).  The types used in the AST are defined in \texttt{ast\_types.mli}, and rely on the OCaml tagged union types to represent the tree.  The OCamlyacc specification for the parser is defined in the \texttt{parser.mly} file.  The parser definition is fairly straight forward.  There is some conversion done at this stage from strings to integers, but primarily it simply builds the AST using the tagged untion types defined in the \texttt{Ast\_types} module.  The root of the generated AST is a \texttt{statement}, specifically a \texttt{Ast\_types.Block} type, which contains a list of \texttt{statement}s.

The \texttt{bawk} compiler can output the AST in a GraphViz format, which is useful for visualizing and debugging the AST.  The \texttt{-ast} flag to the compiler will output GraphViz code.  This can be turned into a PDF by piping it to GraphViz using the command
\begin{center}\tt
bawk .... | dot -Tpdf -o out-file.pdf\rm.
\end{center}
The logic to generate the GraphViz code is the majority of the \texttt{ast.ml} file.  Other functions in that file serve to provide a few helper functions used when parsing the AST.

\bigskip
In the translation phase of the compiler, most of the heavy lifting is done by translating AST structures into basic blocks.  In addition, the compiler keeps track of a a symbol table in the \texttt{env} type in \texttt{compile.ml}.  The structure forms a linked list, with a new node added each time a new, nested scope, is entered.  Each node has a pointer to the parent scope, and a \texttt{StringMap} from symbols to either (1) the jump address for a subroutine, or (2) the index into globals.  The environment also contains additional information for binding variables, which includes the size of the binding.

\subsection{Bytecode operations}
The compiler takes the AST created by the parser portion and translates it to a custom stack-based bytecode.  The instructions for the bytecode are defined in \texttt{bytecode\_types.mli}.  Most of the instructions are related directly to the stack, but several are designed specifically for the application of this language: reading binary files.

The translation phase from the AST to bytecode is really a two phase translation process.  The first does the basic translation from the AST to basic blocks of bytecode.  In this initial translation the various branch instructions are given a unique value referencing a \texttt{Label} instruction.  The \texttt{Label} instruction is a special pseudo instruction that is removed in the second phase of the translation.  When the \texttt{Label} instructions are removed, their address, or position in the instruction list, is noted.  These addresses then replace the values in all the branch instructions, converting them from the value referencing a label to the address of that label.

All the branch instructions in the bytecode operate on absolute addresses.  There is no relative addressing in the bytecode by design, simply for simplicity.  After the second phase of the translation is complete, the \texttt{Label} instructions removed and branch instructions pointing to addresses, the bytecode is ready for the runtime interpreter.

\subsection{Bytecode runtime}
The runtime interpreter is fairly simple, and defined primarily in \texttt{bytecode.ml}.  The runtime executes the stack-based bytecode in a virtual machine environment.  For simplicity the bytecode instructions composing the program, the execution stack, and global variable values are stored in three different arrays: \texttt{instructions}, \texttt{stack}, and \texttt{globals} respectively.  Execution is performed by the recursive \texttt{exec} function, which takes a frame pointer, a stack pointer, and a program counter.

The program counter is simply an index into the \texttt{instructions} array for the instruction being executed in the current cycle.  When the recursive call to \texttt{exec} is made, to execute the next instruction, the program counter is either incremented by one to run the next sequential instruction, or, in the case of a branch or function return instruction, called with some absolute address from the stack or an instruction argument.  Since all branching is to absolute addresses the bytecode interpreter never computes relative offsets from the current program counter.

The stack grows from an index of zero in the \texttt{stack} array upwards, with items pushed and popped as individual instructions are executed.  The argument to \texttt{exec} is a pointer to where the next value can be pushed on the stack.  The stack only holds integer values, which is the primary motivation for not supporting functions, such as a \texttt{printf}-like function, that take string arguments.

Function arguments are placed on the stack, and the frame pointer references the location of these arguments.  The previous frame pointer is also placed onto the stack and reset when the function returns, allowing recursion with arguments to work properly.

Global variables are stored in the \texttt{globals} array, with each symbol, other than function arguments, assigned an offset into this array.  The name `globals' is a bit of a misnomer, since variables local to blocks or functions are also stored in this same array.  The slots in the \texttt{globals} array are not re-used; they last the entire duration of the program and are statically assigned to symbols by the compiler.

The bytecode interpreter \texttt{stack} is limited to 1024 values, and the \texttt{globals} array is limited to this same size of 1024 values.  These limits are \emph{not} enforced by the compiler, but the bytecode interpreter will crash if they are exceeded.  The length of the \texttt{instructions} array is not explicitly limited.


\chapter{Test plan}
The testing for this project was performed using a variety of very tiny bawk language programs designed to exercise various features of the language.  These programs are provided with the source code in the \texttt{tests} directory.  Each program is named \texttt{test-<name>.bawk} and its output is expected to be identical to the \texttt{test-<name>.bawk.out} file.  Both the source and expected output files are listed in \S\ref{sec:test-files}.

The tests are all designed to operate on the binary file \texttt{tests/lichtenstein.png}, which serves as a modestly sized sample file.  The beginning of the \texttt{lichtenstein.png} file is shown below for convenience, and the file itself can be found packaged with the bawk source.
\begin{verbatim}
0000000: 8950 4e47 0d0a 1a0a 0000 000d 4948 4452  .PNG........IHDR
0000010: 0000 0200 0000 0200 0802 0000 007b 1a43  .............{.C
0000020: ad00 0000 0970 4859 7300 000a f000 000a  .....pHYs.......
0000030: f001 42ac 3498 0000 0007 7449 4d45 07d7  ..B.4.....tIME..
0000040: 0511 0921 0919 38af 7500 0020 0049 4441  ...!..8.u.. .IDA
0000050: 5478 dab4 bd5b ac6d 5b76 1dd4 5aef 63ae  Tx...[.m[v..Z.c.
0000060: bdcf bdb7 1ed7 a9b2 2b71 d995 38b1 63f2  ........+q..8.c.
0000070: b021 4fe7 0181 8404 a280 1421 10e2 2522  .!O........!..%"
0000080: 7ef8 46fc 874f 7ef9 e40f 81f8 43e1 2bca  ~.F..O~.....C.+.
0000090: 4f82 2148 0ea0 5809 2838 ca8b d871 1c95  O.!H..X.(8...q..
00000a0: 53b6 eb75 5fe7 ec35 476f 7cf4 3ec6 1c73  S..u_..5Go|.>..s
00000b0: edb5 cfbd 65c8 aeab 5be7 eeb3 f75a 73cd  ....e...[....Zs.
00000c0: 3946 1fbd b7d6 7aeb fcbd fff6 7f2b 7a17  9F....z......+z.
00000d0: 4220 626b 6c06 2a48 9a39 e980 0114 0420  B bkl.*H.9..... 
00000e0: 22f6 d09b c075 8fab 60f4 2638 f317 ed72  "....u..`.&8...r
00000f0: d1e6 7cd8 daa5 b953 4e38 e194 2481 8289  ..|....SN8..$...
...
0059270: 3122 4a04 2150 a6c8 ff1f 00c1 c91a 2379  1"J.!P........#y
0059280: 1a27 0000 0000 4945 4e44 ae42 6082       .'....IEND.B`.
\end{verbatim}

The test programs are executed by a shell script, \texttt{run\_tests.sh}, which runs each program and compares its output to the expected values.


\chapter{Lessons learned}
This project provided a number of valuable software engineering and language design lessons.  The language design lessons were invaluable, as this was my first attempt at writing a compiler from scratch that had any level of sophistication beyond regular expression matching on text files.

\section{Language design}
The design of the language was critical to being successful with this project.  It was crucial to really think through the specifics of the language before even thinking about writing the compiler's translation logic.  It was critical to write the scanner and parser in parallel with hashing out the language details.  Writing these helped force the resolution of many shift/reduce ambiguities and couple reduce/reduce ambiguities in the original, pre-parser, design.

One visible addition to the language that came about from writing the parser is the \texttt{def} keyword before function definitions.  This was added to resolve a conflict in the parser.  Similarly, when adding if statements I realized that parenthesis were required around the conditional expression when a parser conflict was resolved after adding them.

\section{Scanner and parser generators}
As described in \S\ref{sec:bytecode-compiler}, there is some ambiguity in the scanner distinguishing between hexadecimal values and identifier names.  I didn't realize this ambiguity in the initial language design until testing the compiler with a hexadecimal constant beginning with a letter.  This ambiguity was obviously not detected as a shift/reduce or reduce/reduce conflict by the parser generator because it was a collision of definitions in the scanner.  Luckily it was possible to change the language definition slightly and add some logic in the parser to resolve the ambiguity.

\section{Compiler and bytecode design}

\subsection{Return statement problem}
\label{sec:lessons-return}
Late in the project I realized that I did not implement any type of \texttt{return} statement in the langauge.  This was an oversight, but unfortunately the design of the bytecode did not make it easy to add at the last minute. The bytecode instruction \texttt{Rts} returns from a subroutine by popping the arguments to the function off the stack.  Unfortunately this instruction requires a hard-coded constant for the number of instructions to pop.  Without changing the bytecode and its interpreter to support reading this value from the stack it is difficult to return from an arbitrary point in the function.  It could likely still be accomplished without a bytecode interpreter change, but it would require some redesign of the \texttt{translate} function in \texttt{compile.ml} to pass around additional information, namely a \texttt{Label} number for the end of the function.  This change was just too large to risk making in the last week of the semester, so the language stands without a \texttt{return} statement.


\chapter{Source code listing}
\lstset{language=Caml,defaultdialect=[Objective]Caml}

\section{Compiler and bytecode interpreter}
\subsection{Scanner \& parser definitions}
\subsubsection{scanner.mll}
The \texttt{scanner.mll} file contains Ocamllex specifications for the bawk scanner.
\lstinputlisting[caption=scanner.mll]{../scanner.mll}
\subsubsection{parser.mly}
The \texttt{parser.mly} file contains Ocamlyacc specifications for the bawk parser.
\lstinputlisting[caption=parser.mly]{../parser.mly}

\subsection{Interfaces}
\lstinputlisting[caption=parser\_help.mli]{../parser_help.mli}
\lstinputlisting[caption=ast\_types.mli]{../ast_types.mli}
\lstinputlisting[caption=bytecode\_types.mli]{../bytecode_types.mli}
\lstinputlisting[caption=ast.mli]{../ast.mli}
\lstinputlisting[caption=compile.mli]{../compile.mli}
\lstinputlisting[caption=bytecode.mli]{../bytecode.mli}
\lstinputlisting[caption=reader.mli]{../reader.mli}
\lstinputlisting[caption=utile.mli]{../utile.mli}

\subsection{Source Files}
\lstinputlisting[caption=bawk.ml]{../bawk.ml}
\lstinputlisting[caption=parser\_help.ml]{../parser_help.ml}
\lstinputlisting[caption=ast.ml]{../ast.ml}
\lstinputlisting[caption=compile.ml]{../compile.ml}
\lstinputlisting[caption=bytecode.ml]{../bytecode.ml}
\lstinputlisting[caption=reader.ml]{../reader.ml}
\lstinputlisting[caption=utile.ml]{../utile.ml}

\section{Test files}
All test files are designed to be run on the input file \texttt{lichtenstein.png} unless specified otherwise.  The beginning of the \texttt{lichtenstein.png} file is shown below for convenience, and the file itself can be found packaged with the bawk source.

\begin{verbatim}
0000000: 8950 4e47 0d0a 1a0a 0000 000d 4948 4452  .PNG........IHDR
0000010: 0000 0200 0000 0200 0802 0000 007b 1a43  .............{.C
0000020: ad00 0000 0970 4859 7300 000a f000 000a  .....pHYs.......
0000030: f001 42ac 3498 0000 0007 7449 4d45 07d7  ..B.4.....tIME..
0000040: 0511 0921 0919 38af 7500 0020 0049 4441  ...!..8.u.. .IDA
0000050: 5478 dab4 bd5b ac6d 5b76 1dd4 5aef 63ae  Tx...[.m[v..Z.c.
0000060: bdcf bdb7 1ed7 a9b2 2b71 d995 38b1 63f2  ........+q..8.c.
0000070: b021 4fe7 0181 8404 a280 1421 10e2 2522  .!O........!..%"
0000080: 7ef8 46fc 874f 7ef9 e40f 81f8 43e1 2bca  ~.F..O~.....C.+.
0000090: 4f82 2148 0ea0 5809 2838 ca8b d871 1c95  O.!H..X.(8...q..
00000a0: 53b6 eb75 5fe7 ec35 476f 7cf4 3ec6 1c73  S..u_..5Go|.>..s
00000b0: edb5 cfbd 65c8 aeab 5be7 eeb3 f75a 73cd  ....e...[....Zs.
00000c0: 3946 1fbd b7d6 7aeb fcbd fff6 7f2b 7a17  9F....z......+z.
00000d0: 4220 626b 6c06 2a48 9a39 e980 0114 0420  B bkl.*H.9..... 
00000e0: 22f6 d09b c075 8fab 60f4 2638 f317 ed72  "....u..`.&8...r
00000f0: d1e6 7cd8 daa5 b953 4e38 e194 2481 8289  ..|....SN8..$...
...
0059270: 3122 4a04 2150 a6c8 ff1f 00c1 c91a 2379  1"J.!P........#y
0059280: 1a27 0000 0000 4945 4e44 ae42 6082       .'....IEND.B`.
\end{verbatim}

\testfile{test-arith.bawk}
\testfile{test-cond1.bawk}
\testfile{test-func1.bawk}
\testfile{test-func2.bawk}
\testfile{test-func3.bawk}
\testfile{test-pat1.bawk}
\testfile{test-pat2.bawk}
\testfile{test-pat3.bawk}
\testfile{test-pat4.bawk}
\testfile{test-pat5.bawk}
\testfile{test-pat6.bawk}
\testfile{test-pat7.bawk}
\testfile{test-pat8.bawk}
\testfile{test-pat9.bawk}
\testfile{test-pat10.bawk}
\testfile{test-scope1.bawk}
\testfile{test-useful1.bawk}
\testfile{test-var1.bawk}

\section{Misc. Files}
\lstinputlisting[caption=Makefile,language=make]{../Makefile}
\lstinputlisting[caption=run\_tests.sh,language=bash]{../run_tests.sh}


\end{document}
